\documentclass[preprint,12pt]{elsarticle}
\input{jmacros.tex}
\journal{SoftwareX}
\hypersetup{pdftitle={SWAFL Desktop: equation-traceable engines for irrigation scheduling and design},pdfauthor={Naziru Halilu}}
\begin{document}

\begin{frontmatter}
\title{SWAFL Desktop: equation-traceable engines for irrigation scheduling and design}

\author[aua,upna]{Naziru Halilu\corref{cor1}}
\ead{halilunaziru73@gmail.com}
\cortext[cor1]{Corresponding author. ORCID: 0009-0006-2324-9599.}
\affiliation[aua]{organization={Agricultural University of Athens},city={Athens},country={Greece}}
\affiliation[upna]{organization={Public University of Navarre},city={Pamplona},country={Spain}}

\begin{abstract}
Irrigation software seldom shows which relation produced a number or where that relation holds. SWAFL Desktop is an open, offline application that couples soil, water, crop and atmosphere for irrigation scheduling and design through three computing engines: a meta-model of a 10\,190-run Daisy simulation ensemble, a daily layered soil-water model, and closed-form surface, drip and sprinkler design modules. Every output is traced to one of 134 numbered equations carrying a provenance label (literature, data-driven or engineered), and every engine is tested against a reference outside the software: published worked examples, held-out simulations, mass closure and an independent solver. The tests reproduce textbook examples within 2\,\%, rank held-out scenarios correctly, and exposed three defects that were corrected.
\end{abstract}

\begin{keyword}
irrigation decision support \sep irrigation design \sep simulation meta-model \sep equation traceability \sep software verification \sep Daisy model
\end{keyword}
\end{frontmatter}

\section*{Code metadata}
\begin{table}[H]\centering\small
\begin{tabularx}{\textwidth}{lp{5.2cm}X}
\toprule
Nr. & Code metadata description & Metadata\\\midrule
C1 & Current code version & v1.1.0\\
C2 & Permanent link to code/repository used for this code version & \url{\repourl}\\
C3 & Permanent link to Reproducible Capsule & Not available; all checks re-run with one command (\texttt{Manuscripts/reproducibility})\\
C4 & Legal Code License & MIT\\
C5 & Code versioning system used & git\\
C6 & Software code languages, tools, and services used & JavaScript (Electron, Node.js), HTML/CSS; Python for the verification scripts\\
C7 & Compilation requirements, operating environments \& dependencies & Node.js and Electron (build); Leaflet and geotiff.js (bundled); runs on Windows, macOS and Linux and in a web browser\\
C8 & If available, link to developer documentation/manual & \url{\manualurl}\\
C9 & Support email for questions & halilunaziru73@gmail.com\\\bottomrule
\end{tabularx}
\end{table}

\section*{Software metadata}
\begin{table}[H]\centering\small
\begin{tabularx}{\textwidth}{lp{5.2cm}X}
\toprule
Nr. & Software metadata description & Metadata\\\midrule
S1 & Current software version & 1.1.0\\
S2 & Permanent link to executables of this version & \url{\exeurl}\\
S3 & Permanent link to Reproducible Capsule & Not available\\
S4 & Legal Software License & MIT (source); end-user licence for the compiled application\\
S5 & Computing platforms/Operating Systems & Microsoft Windows (portable executable); macOS and Linux (Electron); any modern web browser\\
S6 & Installation requirements \& dependencies & None for the Windows executable; Node.js for running from source\\
S7 & If available, link to user manual & \url{\manualurl}\\
S8 & Support email for questions & halilunaziru73@gmail.com\\\bottomrule
\end{tabularx}
\end{table}

\linenumbers
\section{Motivation and significance}
Between soil-water science and a farmer's decision about when, how much and by which method to irrigate lies a layer of software: schedulers, sensor dashboards, design calculators and crop-water models~\cite{Datta2023,Garcia2020,Jones2004}. Three weaknesses recur in that layer. Tools report numbers without saying which equation produced them or whether the equation comes from the literature or from an in-house heuristic; they are tested on cases chosen by their authors, with results seldom published in a form a third party can re-run~\cite{Hutton2016,Wilkinson2016}; and screening emulators are sometimes used as absolute predictors~\cite{Bennett2013,Moriasi2007}. Process models such as Daisy~\cite{Hansen1991,Hansen2012,Holbak2026}, HYDRUS~\cite{Simunek2008}, DSSAT~\cite{Jones2003}, APSIM~\cite{Keating2003} and AquaCrop~\cite{Steduto2009} are the scientific backbone, but they are not decision tools for a user who needs an answer in seconds. Emulators bridge that gap~\cite{Razavi2012,Asher2015}, yet their literature concerns statistical accuracy more than traceability. Meanwhile, classical irrigation design (volume-balance advance in surface irrigation~\cite{Walker1987,Walker1989}, lateral hydraulics in drip systems~\cite{KellerKarmeli1974,Keller1990}, nozzle and pattern analysis for sprinklers~\cite{Keller1990,Christiansen1942,Playan2005}) is published as closed-form equations with worked examples, which are a free and under-used source of verification data.

SWAFL Desktop addresses these weaknesses under three design principles: every number is equation-traceable, every equation carries a provenance label, and every claim of fidelity is backed by a re-runnable test against a reference from outside the software. The application is intended for irrigation engineers, agronomists, extension services, teachers and students who need transparent first estimates of irrigation schedules and system designs, and for researchers who want to inspect how a decision depends on soil, climate and management.

\section{Software description}
\subsection{Software architecture}
SWAFL Desktop is an Electron application written in JavaScript that runs offline on Windows, macOS and Linux and in a web browser. All computation runs locally; the simulation tables are embedded as JSON. \Cref{fig:arch} summarises the architecture: an evidence layer (the archived simulation ensemble), two engines, three design layers, an equation-traceable dashboard and the decision outputs.

\emph{Engine~1} is an annual meta-model. It composes, by multiplication of ratios and addition of differences, 20-year-mean results of 10\,190 Daisy simulations in 17 designed experiments~\cite{HaliluData} to give evapotranspiration, soil evaporation, transpiration, irrigation, yield, drainage and nitrogen leaching for the chosen soil, tillage, mulch, rainfall scale, warming, trigger, sensor depth, irrigation method and nitrogen rate. It fits no parameter, so it cannot leave the simulated envelope without clipping.

\emph{Engine~2} is a daily 40-layer bucket model over 20 years of weather. It produces dates, depths, the animated field water balance and the irrigation schedule, and is scaled so that its 20-year means equal those of Engine~1, so that the dates and the annual numbers describe the same soil.

The \emph{design layers} use hydraulics rather than data: the surface layer (border, furrow, basin) solves the volume-balance advance with Kostiakov--Lewis intake~\cite{Walker1987,Walker1989}, the drip layer solves the lateral by iterating emitter discharge and Hazen--Williams friction to convergence~\cite{KellerKarmeli1974,Keller1990}, and the sprinkler layer applies orifice discharge, range, Christiansen uniformity and the spacing rules of the textbook~\cite{Keller1990,Christiansen1942,Playan2005}. They take soil, crop and trigger from the engines.

%%PULL p2.tex fig:arch side

\subsection{Software functionalities}
The user works in one window with a tree of layers, the drawing of the layer, a dashboard and a panel of project properties (soil, tillage, crop, mulch, weather, irrigation method, trigger, nitrogen and the working date from 1980 to 2100). The main functions are:
\begin{itemize}[nosep]
\item \textbf{Irrigation scheduling}: sensor depth and suction trigger against 20 years of daily soil water; the trigger law, which expresses the triggers of ten soils as a common root-zone depletion; next irrigation, trigger date and wilting date for the day on the clock.
\item \textbf{Irrigation design}: layout, advance and recession, infiltrated depth, efficiencies and a design table for surface irrigation; lateral hydraulics, emission uniformity and wetted bulb for drip; nozzle, spacing, uniformity, pipe network and hydraulic grade line for sprinklers; a search for the most efficient inputs.
\item \textbf{Soil, mulch and climate exploration}: the fate of the water a mulch saves, evaporation regimes, tillage and consolidation, root structure, seven soil--atmosphere feedback loops, and 20-season reliability, water and nitrogen budgets.
\item \textbf{Traceability}: each tile, chart and arrow names its equation; the 1\,154-page manual~\cite{HaliluManual} prints all 134 numbered equations of the engines with their provenance (38 literature, 18 data-driven, 78 engineered), and the code cites the equation numbers.
\item \textbf{Reproducibility}: one command re-runs every verification test; results, figures and reports are exported as files.
\end{itemize}
Single-call medians in Node.js are 0.021\,ms for Engine~1, 0.25\,ms for the drip layer, 1.7\,ms for the sprinkler layer and 0.14\,s for a 20-year Engine~2 simulation, so a sweep of $10^5$ Engine~1 scenarios takes about 2\,s and interaction is limited by rendering, not computation.

\section{Illustrative examples}
\subsection{Surface irrigation against a textbook}
Examples~6.1 (furrow) and 6.4 (border) of Walker and Skogerboe \cite{Walker1987} were recomputed. For the furrow, 16 published quantities were reproduced with a median error of 0.08\,\% and a maximum of 1.81\,\% (advance time at $Q=0.10$\,m$^3$\,min$^{-1}$: 81.65 against 80.2\,min), and the required infiltration times agreed to better than 0.01\,\% (\cref{fig:ws}). All border quantities were within 0.5\,\% (\cref{tab:border}). Thirteen advance times of the FAO Irrigation Paper~45 border examples~\cite{Walker1989} were reproduced within $-0.2$ to $+0.5$\,\%. The test also exposed a defect: release~1.0.0 used a fixed advance exponent in the steady-intake term, which made slow advance too fast (up to $-29$\,\%); iterating the exponent, as the textbook does, brought every case within $-0.2$ to $+1.8$\,\% in release~1.1.0.

%%PULL p2.tex fig:ws
%%PULL p2.tex tab:border

\subsection{Drip and sprinkler layers}
The drip lateral loss agrees with the Christiansen reduction factor and the exact discrete sum to 0.01\,\% or better for 5 to 400 outlets. Against an independent solver on 400 random configurations, the first version (four iterations) deviated by up to 81\,\% where the head loss exceeded 20\,\% of the inlet head; the rebuilt solver iterates to convergence, agrees to $3.2\times10^{-7}$ for $h_f/H_0\le0.6$, and warns the user beyond 0.2 (\cref{fig:drip}). For the sprinkler textbook example the software follows the book's decision logic (number of sets, laterals and days), the nozzle discharge is within 1.3\,\%, and the wetted diameter, which the first version under-predicted by 14\,\% until its range coefficient was replaced by the published value, is now within 1\,\%. One limitation remains visible to the user: the engineered wind term changes uniformity by less than 0.2 percentage points at 3\,m\,s$^{-1}$ and should not be used for design.

%%PULL p2.tex fig:drip

\subsection{Engines against held-out simulations and mass balance}
Engine~1 was queried at the factor levels of 1\,000 held-out Daisy runs (sensor depth $\times$ trigger $\times$ rain on a coarse loam). It ranked them correctly (Spearman $\rho=0.96$ for irrigation, 0.92 for drainage, 0.97 for soil evaporation), but its raw values carried a baseline offset between experiment families ($+53$\,mm for evapotranspiration, $+71$\,mm for irrigation); after aligning the baseline, the normalised RMSE was 6.2\,\% for evapotranspiration and 18.9\,\% for irrigation (\cref{fig:e1}, \cref{tab:e1}). Engine~1 is therefore a rank-faithful screening emulator, and its yield output is for ranking only. The unscaled daily core of Engine~2 closed the water balance to a median 0.1\,\% of inflow, and on a grid of 10 soils $\times$ 5 triggers $\times$ 2 methods its annual irrigation agreed with Engine~1 within 10\,\% in 90 of 100 scenarios; the ten exceptions are coarse soils at dry triggers, where the bucket never reaches the trigger (\cref{fig:e2}).

%%PULL p2.tex fig:e1
%%PULL p2.tex tab:e1
%%PULL p2.tex fig:e2

\section{Impact}
SWAFL Desktop turns a large, published process-model ensemble and the classical design theory of three irrigation methods into an interactive, offline tool whose every number can be audited: a reviewer can select a displayed quantity, find its equation in the manual, check the equation against its source and re-run the test that supports it, without access to the developer. This changes the questions that can be asked of decision software. Instead of ``is the tool right?'', which a field trial must answer, one can ask which outputs rest on published theory reproduced against its textbook (the surface and drip layers), which on simulated ratios that preserve the ranking of options (Engine~1), and which on engineered shapes that should be confirmed with manufacturer data (sprinkler patterns, wind drift). The verification described here found three defects that were invisible on the screen, which supports the practice of shipping a published-example test with every closed-form layer.

The software supports research on soil--water--atmosphere interactions (for example the drainage share of mulch savings and the reach of suction triggers on different soils), teaching of irrigation engineering, where a student can follow a number from the dashboard to the equation and the test, and extension work in places without reliable internet, because the single Windows executable needs no installation and no network. The source, the manual, the verification scripts and the simulation data are open, so the evaluation can be repeated and extended by any group. No field validation, agronomic benefit or user study is claimed; the operating envelopes reported here define where field measurements would be most informative.

\section{Conclusions}
SWAFL Desktop 1.1.0 is an open, offline application for irrigation scheduling and design in which every output is tied to a numbered equation, a provenance label and a re-runnable test. Its design layers reproduce published worked examples within 2\,\%, its annual meta-model ranks held-out scenarios correctly with a known baseline offset, its daily model conserves mass, and its drip solver is reliable up to a head loss of 60\,\% of the inlet head with a warning beyond 20\,\%. Verification exposed three defects, all corrected in the current release; the remaining open points (the baseline offset, three unverified sprinkler range coefficients, the wind term) are stated in the software and the manual. Field validation is the next step.

\section*{CRediT authorship contribution statement}
\textbf{Naziru Halilu:} Conceptualization, Methodology, Software, Validation, Formal analysis, Data curation, Visualization, Writing -- original draft, Writing -- review \& editing.

\section*{Declaration of competing interest}
The author declares that he has no known competing financial interests or personal relationships that could have appeared to influence the work reported in this paper.

\section*{Declaration of generative AI and AI-assisted technologies in the writing process}
During the preparation of this work the author used Claude (Anthropic) to edit language and to format the manuscript for the journal. After using this tool, the author reviewed and edited the content as needed and takes full responsibility for the content of the publication.

\section*{Funding}
This research did not receive any specific grant from funding agencies in the public, commercial, or not-for-profit sectors.

\section*{Data availability}
The simulation ensemble behind Engine~1 is archived on Zenodo (\doihl{10.5281/zenodo.23043626}); every log of every run is in the repository folder \cslink{}. The verification scripts are in \url{\codeurl}.

\section*{Acknowledgements}
The Daisy model is developed and maintained at the University of Copenhagen.

\bibliographystyle{elsarticle-num}
\bibliography{../tex/refs}
\end{document}
